% rosetta/appendix-table.tex — 附录 A：一行一对的总表（由 tools/make_table.py 生成，勿手改）
\section{对照总表：一行一对}
\label{app:table}

本表把正文 107 条词条压成一行一对（第一列词条名即「数学侧」，第二列为其英文原词，第三列取该词条 \texttt{bridge}（对应）段首句的骨架，长公式从略），供快速检索；定义、例子、直觉与易错点见正文各轨道词条。

另有 Nahm 方程、Nahm 定理、相干层、导出范畴、环路群、opers、Dolbeault 上同调、Kac--Moody 代数 等 8 条原文只引用而未实质使用的背景术语（\S 9.3 标 $\dagger$），其最小可用定义见附录 B。

\begingroup
\footnotesize
\renewcommand{\arraystretch}{1.12}
\begin{longtable}{@{}p{4.3cm}p{2.6cm}p{8.8cm}@{}}
\toprule
\bfseries 词条（数学侧） & \bfseries 英文原词 & \bfseries 一句话对应 \\
\midrule
\endfirsthead
\toprule
\bfseries 词条（数学侧） & \bfseries 英文原词 & \bfseries 一句话对应 \\
\midrule
\endhead
\bottomrule
\endlastfoot

\multicolumn{3}{@{}l}{\bfseries T1 对偶、群论与耦合常数}\\
\addlinespace[2pt]
电-磁对偶 & electric--magnetic duality & §1 p.2--3：物理侧的电-磁对偶群就是根数据对偶化所得的 ${}^LG$，耦合按…… \\
Montonen--Olive 对偶 & Montonen--Olive duality & §1 p.2：KW06 把 Montonen--Olive 猜想当作出发点，真正的工具是它的精确版本 $S$-对偶（§2.2） \\
S-对偶 & S-duality & §2.2 p.11：$S$ 的矩阵为 （单连通型）或 \\
T-对偶 & T-duality & §1 p.3：Hitchin 模空间的 $T$-对偶在数学上被理解为 Fourier--Mukai 变换，是几何 Langlands 的某种…… \\
GNO 双群 & GNO dual group & §1 p.3：KW06 明确说「对规范群 $G$，GNO 双群实际上就是 Langlands 对偶群 ${}^LG$」，并归功于 Atiya…… \\
Langlands 对偶群 & Langlands dual group & §1 p.3--4 与 §2.2 p.11：数学的 ${}^LG$ 与物理的对偶群是同一个构造 \\
单连通型 & simply-laced & §2.2 p.11：原文先写单连通型时 $S= $ 并生成 $SL(2,\mathbb Z)$，再写一般情形的 $S= $ \\
对偶群 $\Gamma$ & duality group & §2.2 p.11：$\Gamma$ 由 $S,T$ 生成，是 $SL(2,\mathbb R)$ 的无限离散子群 \\
$\theta$ 角 & theta angle & §1 p.3：$\theta$ 与 $e$ 一起给出 $\tau$ \\
复耦合 $\tau$ & complexified coupling & §2.2 p.11：$\tau=\theta/2\pi+4\pi i/e^2$，$S$ 与 $T$ 生成 $\Gamma$ \\
权格 & weight lattice & §2.2 p.13：$\vec n$（电）取值于 $\mathfrak g$ 的权格、$\vec m$（磁）取值于余权格 \\
余权格 & coweight lattice & §2.2 p.13：$\vec m$ 取值于余权格 \\
GNO 线 & GNO line & §6.2 p.84：「$G$ 的 't Hooft 算子由 ${}^LG$ 的表示分类」就是经典 GNO 对偶，近来被重解为「关于算子而非态…… \\
Hecke 群 & Hecke group & §2.2 p.11：非单连通型时 $S:\tau\mapsto-1/(n_{\mathfrak g}\tau)$，生成的是 Hecke 群而…… \\
\multicolumn{3}{@{}l}{\bfseries T2 超对称、扭化与拓扑场论}\\
\addlinespace[2pt]
$N=4$ 超对称杨--米尔斯 & $N=4$ supersymmetric Yang--Mills & 物理侧这一个理论按 §3.1 的三种方式扭化，给出三种拓扑场论 \\
超荷 & supercharge & 扭化直接作用在超荷上：$\epsilon\in\mathcal S^+$ 在…… \\
标量超荷 & scalar supercharge & 标量超荷是「物理侧能翻译成数学」的入口：$Q$ 是微分，$Q$-上同调是上同调 \\
BRST 算子 & BRST operator & 「BRST 算子 = 标量超荷 = 拓扑超荷」是同一对象在三种语境下的名字：数学侧叫微分，物理侧叫超荷，规范理论传统叫 BRST \\
$Q$-上同调 & cohomology of $Q$ & $Q$-上同调是词典的翻译机制：物理侧的 $Q$-上同调类对应数学侧的上同调类 \\
R-对称 & R-symmetry & 扭化的定义就是取同态 ，用 $Spin'(4)=(1\times\varkappa)(Spin(4))$ 替换原转动群 \\
扭化 & twisting & 扭化是两个独立事实的合成：其一，超荷按 R-对称的某个子群（这里是 ）重新定义自旋，从而让标量超荷存在 \\
GL 扭化 & GL twist & GL 扭化是几何 Langlands 的物理载体：实 $t$ 侧是 A-模型（靶 $\mathcal M_H$ 的辛结构）…… \\
拓扑场论 & topological field theory & 拓扑场论是扭化的输出、也是词典的靶：物理侧的 TQFT 对应数学侧的不变量理论（Donaldson 不变量、Gromov--Witten 不…… \\
瞬子数 & instanton number & 瞬子数把物理侧的拓扑荷翻译成数学侧的特征数：$k=\langle c_2(E),[M]\rangle$ \\
BPS 条件 & BPS condition & BPS 条件把物理的「质量等于中心荷」翻译成数学的「稳定/自对偶方程」：Hitchin 方程、单极子方程、瞬子方程都是 BPS 方程 \\
中心荷 & central charge & 中心荷是 BPS 质量的物理来源，也是数学侧周期积分的物理对应：物理的电荷格对应数学的 $H_1(\Sigma,\mathbb Z)$，物理…… \\
拓扑超荷 & topological supercharge & 拓扑超荷是词典的「动词」：它把物理侧的对称性变成数学侧的微分，把物理的不变量变成数学的上同调类 \\
Coulomb 分支 & Coulomb branch & Coulomb 分支是「物理侧能算」的地方：中心荷、BPS 质量、$S$-对偶作用都在这里显式计算，再借超对称与全纯性延拓到整个理论（§2.…… \\
\multicolumn{3}{@{}l}{\bfseries T3 Hitchin 模空间与超凯勒几何}\\
\addlinespace[2pt]
Hitchin 模空间 & Hitchin moduli space & 数学侧的 $\mathcal M_H$ 与物理侧的低能靶空间是同一个实流形：它的点在复结构 $I$ 下是半稳定 Higgs 丛、在复结构…… \\
Higgs 丛 & Higgs bundle & $\mathcal M_H$ 在复结构 $I$ 下恰是半稳定 Higgs 丛的模空间（含严格半稳定的某些等价类），即原文 §4.2 的操作…… \\
Higgs 场 & Higgs field & 物理侧的 $\phi$ 与数学侧的 $\varphi$ 是同一个场在两个复结构下的两种读法：在复结构 $I$ 中读 $A_{\bar z}$…… \\
Hitchin 对 & Hitchin pair & $I$ 结构下 $\mathcal M_H$ 等同于稳定 Hitchin 对的模空间（Hitchin 1987，原文 §4.2）：全纯 mo…… \\
稳定性 & stability & $\mathcal M_H$ 只参数化稳定 Hitchin 对与严格半稳定对的某些等价类（原文 §4.2）：这正是操作 (II)$'$ 与操…… \\
超凯勒流形 & hyper-Kähler manifold & 数学侧的超凯勒结构就是物理侧 $(4,4)$ 超对称的几何翻译：三个复结构 $I,J,K$ 对应三个超荷方向，三个辛形式…… \\
超凯勒商 & hyper-Kähler quotient & 超凯勒商是「Hitchin 方程解模规范变换」的精确几何表述：$\vec\mu=0$ 就是 Hitchin 方程的三个实方程，除以…… \\
复结构三元组 & complex structures $I,J,K$ & $S$-对偶/镜像对称交换的是复结构 $I$ 与 $J$ 两侧（A-模型与 B-模型靶空间互换）：$B$-模型 on $Y({}^LG)$…… \\
复结构 $\mathbb P^1_g$ & twistor $mathbb P^1$ & twistor $\mathbb P^1$ 与四维 twist 参数空间的对应 $w_+=-t$、$w_-=t^{-1}$ 把「三个复结构」…… \\
辛形式 & symplectic form & 物理侧的 $\theta$ 项、$B$-场与 A-模型的辛结构都由数学侧的 $\omega_I,\omega_J,\omega_K$ 给出…… \\
moment map $\mu$ & moment map & 一一对应：三个 moment map 同时为零 ⟺ Hitchin 方程 \\
Hitchin 纤维化 & Hitchin fibration & 数学侧：$\pi$ 把 Higgs 丛映到特征多项式，是代数完全可积系统的动量映射 \\
可积系统 & integrable system & $B$ 的维数恰为 $\mathcal M_H$ 的一半，且纤维是 Lagrange 环面，这就是「可积」的全部内容 \\
Casimir 基（Hitchin 基 $B$） & Casimir basis & 数学侧：$B$ 是 Hitchin 纤维化的基，也是 Chevalley 定理…… \\
谱覆盖 & spectral cover & 数学侧：谱覆盖给出 Hitchin 纤维的具体模型（Jacobi/Prym 簇），是可积系统显式化的标准技术 \\
\multicolumn{3}{@{}l}{\bfseries T4 广义复几何、sigma 模型与膜}\\
\addlinespace[2pt]
sigma 模型 & sigma model & 数学侧给 $X$ 一个几何结构（复结构、辛结构、超凯勒三元组），物理侧就得到一个拓扑 sigma 模型 \\
A-模型 & A-model & A-模型对应辛结构，B-模型对应复结构 \\
B-模型 & B-model & 对象级：层 $\mathcal F$ 映到支撑在 $\mathrm{supp}\,\mathcal F$ 上、Chan-Paton 丛为…… \\
镜像对称 & mirror symmetry & 镜像对称等于在 Hitchin 纤维的环面上作 $T$-对偶、同时在基上加一个映射 $\Xi$ \\
广义复几何 & generalized complex geometry & 广义复几何是 A/B 模型共居的舞台：$\mathcal I_J$ 给 B-模型，$\mathcal I_\omega$ 给 A-模型…… \\
广义复结构 & generalized complex structure & 广义复结构是 A-模型与 B-模型的「共同母体」：$\mathcal I_J$ 与 $\mathcal I_\omega$ 分别对应复结构与…… \\
B-场 & B-field & $B$-场在数学侧有三个化身：复化 Kähler 类 $[B+i\omega]$（A-模型的参数）、广义复结构的 B-变换…… \\
余各向同性 & coisotropic & 余各向同性是「辛意义下的子对象」：它由 Poisson 对易关系（第一类约束）定义，因而在数学上对应 Poisson 几何的余各向同性子流形…… \\
A-膜 & A-brane & A-膜对应辛几何对象 \\
B-膜 & B-brane & 对象级：相干层 $\mathcal F$ 映到支撑在 $\mathrm{supp}\,\mathcal F$ 上、Chan-Paton 为…… \\
膜 & brane & 膜对应范畴的对象 \\
边界条件 & boundary condition & 边界条件等同于膜，等同于范畴的对象 \\
Chan-Paton 丛 & Chan-Paton bundle & Chan-Paton 丛就是膜的「层」那一半数据：数学侧的对象是（支撑，层），其不变量是层上同调 \\
零膜 & zerobrane & 零膜对应点层 $\mathcal O_p$（$B$-模型侧） \\
type F 膜 & type F brane & type F 膜对应「Hitchin 纤维加平坦 $U(1)$ 丛」，也就是谱曲线/Prym 簇上的一个点 \\
典范余各向同性 A-膜 & canonical coisotropic A-brane & c.c. 膜是把 A-膜翻译成 $D$-模的机器：函子 $A\mapsto\mathrm{Hom}(A_{cc},A)$ 把…… \\
膜的类型 & brane type & 三个字母对应三个复结构 $I,J,K$，不是三种膜、也不是三个模型 \\
SYZ 图景 & SYZ picture & SYZ 图景对应镜像对称的几何实现，也对应 Fourier--Mukai 变换：纤维上的 $T$-对偶在 $B$-膜范畴里表现为核层给出的…… \\
\multicolumn{3}{@{}l}{\bfseries T5 线算子、特征膜、单极子与 Hecke}\\
\addlinespace[2pt]
Wilson 算子 & Wilson operator & 对应规则：Wilson 算子 $W_R(S)$ 对应 Hecke 算子 $H_{R,x}$，其中 $x$ 是 $S$ 与黎曼面 $C$ 的交…… \\
't Hooft 算子 & 't Hooft operator & 核心对应：${}^LG$ 的最高权 ${}^Lw$（GNO 磁荷）$\leftrightarrow$ 仿射 Grassmann 流形中的点…… \\
线算子 & line operator & 对应规则：线算子的作用范畴与 Hecke 算子的作用范畴同构 \\
圈算子 & loop operator & 对应规则：圈算子对应 Hecke 特征层的「特征值」一侧——在闭圈上取迹相当于在 Hecke 特征层上取 $V_E$ 的纤维维数 \\
有序算子 & order operator & 对应规则：有序算子对应「张量积型」的 Hecke 作用（Wilson 侧），无序算子对应「修正型」的 Hecke 作用（'t Hooft 侧…… \\
无序算子 & disorder operator & 对应规则：无序算子 $\leftrightarrow$ Hecke 修正（改写丛），其标签是 ${}^LG$ 的最高权 \\
特征膜 & eigenbrane & 对应规则：特征膜 $\leftrightarrow$ Hecke 特征层 \\
电特征膜 & electric eigenbrane & 对应规则：电特征膜 $\leftrightarrow$ Hecke 特征层（$\mathrm{Bun}_G$ 上的 $D$-模） \\
磁特征膜 & magnetic eigenbrane & 对应规则：磁特征膜 $\leftrightarrow$ Hecke 特征层在 S-对偶下的像，其标签是 ${}^LG$ 的表示…… \\
Hecke 特征层 & Hecke eigensheaf & 对应规则：Hecke 特征层 $\leftrightarrow$ 特征膜 \\
Hecke 算子 & Hecke operator & 对应规则：Hecke 算子 $H_{V,x} \leftrightarrow$ 't Hooft 线算子 $T(V,x)$（$\Psi=0$…… \\
Hecke 修正 & Hecke modification & 对应规则：修正的型 $\leftrightarrow$ 't Hooft 算子的 GNO 磁荷…… \\
Hecke 对应 & Hecke correspondence & 对应规则：Hecke 对应 $\leftrightarrow$ 't Hooft 算子的时空轨迹…… \\
单极子 & monopole & 对应规则：单极子 $\leftrightarrow$ Hecke 算子的几何实现 \\
单极子冒泡 & monopole bubbling & 对应规则：冒泡 $\leftrightarrow$ associated 权（Schubert 循环的边界层），即 Hecke 修正空间紧化…… \\
Bogomolny 方程 & Bogomolny equations & 对应规则：Bogomolny 方程的解空间 $\leftrightarrow$ Hecke 对应 $Q$ 的纤维…… \\
仿射 Grassmann 流形 & affine Grassmannian & 对应规则：$\mathrm{Gr}_G$ 的点 $\leftrightarrow$ Hecke 修正的数据…… \\
Schubert 胞腔 & Schubert cell & 对应规则：Schubert 胞腔 $\leftrightarrow$ 固定磁荷的修正空间 $\leftrightarrow$ 一个 Heck…… \\
Schubert 循环 & Schubert cycle & 对应规则：Schubert 循环 $\leftrightarrow$ associated 权型 $\leftrightarrow$ 修正空…… \\
泛丛 & universal bundle & 对应规则：泛丛 $\mathcal{E} \leftrightarrow$ Hecke 特征层公式中的…… \\
超多重态 & hypermultiplet & 对应规则：超多重态 $\leftrightarrow$ Higgs 场 $\varphi$（其四个实标量在扭化后组合成复 1-形式…… \\
\multicolumn{3}{@{}l}{\bfseries T6 Langlands 对应、$\mathcal D$-模与几何背景}\\
\addlinespace[2pt]
几何 Langlands & Geometric Langlands & 两边的对象：$A$-膜 $\leftrightarrow \mathrm{Bun}_G(C)$ 上的扭化 ${\cal D}$-模 \\
Langlands 纲领 & Langlands Program & 对应：GNO 把电通量格 $X_*(T)$ 与磁通量格 $X^*(T)$ 互换，正是 Langlands 对偶交换 weight/cowei…… \\
D-模 & D-module & 对应：符号分次 给出余切丛上的函数代数，故 ${\cal D}_X$-模是「量子化余切丛」上的对象 \\
微分算子层 & Sheaf of Differential Operators & 对应：符号映射 $\sigma:{\cal D}_X\to\mathrm{Sym}\,T_X$ 就是经典--量子对应 \\
扭化 D-模 & Twisted D-module & 对应：扭化类在物理侧是背景通量，在数学侧是线丛的转移数据 \\
平坦联络 & Flat Connection & 对应：平坦联络 $=$ 局部系统 $= \pi_1(C)$ 的表示 \\
局部系统 & Local System & 对应：$\mathrm{Loc}_{^LG}$ 是特征簇的 Betti 版本，平坦联络模空间 $\mathcal Y(^LG,C)$ 是 d…… \\
基本群 & Fundamental Group & 对应：$\pi_1(C)$ 到 $^LG$ 的同态 $\leftrightarrow C$ 上平坦 $^LG$-丛…… \\
不可约表示 & Irreducible Representation & 对应：$E$ 不可约 $\leftrightarrow F_E$ 唯一且不可约 \\
分歧 / 非分歧 & Ramified / Unramified & 对应：孔 $\leftrightarrow$ Hecke 算子的插入点 \\
主丛 & Principal Bundle & 对应：主丛 $=$ 规范场论的拓扑背景 \\
向量丛 & Vector Bundle & 对应：主 $G$-丛加表示 $\rho$ 给出向量丛 $E\times_G V$ \\
模空间 & Moduli Space & 对应：$\mathcal M(G,C)$ 上的扭化 ${\cal D}$-模…… \\
Casimir 多项式 & Casimir Polynomial & 对应：Casimir 多项式 $\leftrightarrow$ 谱覆盖系数 $\leftrightarrow$ 纤维上阿贝尔簇的参数 \\
量子化 & Quantization & 对应：${\cal D}$-模是「余切丛量子化」在对象层面的体现，「$A$-膜 $\to {\cal D}$-模」这一步就是量子化翻译 \\
非交换几何 & Noncommutative Geometry & 对应：${\cal D}$-模范畴 $=$ 余切丛非交换形变上的 $B$-膜范畴 \\
Fourier--Mukai 变换 & Fourier--Mukai Transform & 对应：FM 把结构层/点送到向量丛（skyscraper 到平坦丛），正是「局部系统 $\leftrightarrow {\cal D}$-…… \\
Kähler 商 & Kähler Quotient & 对应：D-项方程 $\leftrightarrow$ moment map \\
通量 & Flux & 对应：电通量格 $X_*(T)$ 与磁通量格 $X^*(T)$ 在 GNO 下互换，正是 $^LG$ 交换 weight/coweight…… \\
de Rham 微分 & de Rham Differential & 对应：${\cal D}$-模的联络 $\nabla$ 是 $d$ 的「丛版本」…… \\
消灭定理 & Vanishing Theorem & 对应：$H^{>0}(C,\mathrm{ad}E)=0$ 这类条件 $\leftrightarrow$ 无穷小形变无阻、真空孤立 \\
\multicolumn{3}{@{}l}{\bfseries T7 深水区：盲测指出的四处机制补篇}\\
\addlinespace[2pt]
单极子模空间与 Hecke 修正的机制 & monopole moduli and Hecke modifications & 对应规则：算子 $T({}^Lw,z_0)$ 对应型 ${}^Lw$ 在 $z_0$ 的修正 \\
Nahm 极点与 ADHMN 构造 & Nahm poles and the ADHMN construction & 对应规则：Nahm 数据 $\leftrightarrow$ 单极子的谱数据 \\
膜到 D-模的翻译器 & the brane-to-D-module functor & 对应规则：$K$ 复结构中的 A-膜范畴 $\to \mathcal D_{K_M^{1/2}}$-模范畴，对象…… \\
三种拓扑扭化 & the three topological twists & 对应规则：三类扭化 $\leftrightarrow$ 保住的标量超荷数（一个、两个无族、两个成族） \\
\end{longtable}
\endgroup
